% GENERATED FILE. Edit canonical lessons, then run npm run generate:books.
% canonical-source-hash: fnv1a64:50af4795be77a8fb
% canonical-lessons: TA-C230-karpanai-cey, TA-C230-cantekappatu, TA-C230-kurai-col, TA-C230-pativirakku, TA-C230-pativerru

\chapter{To imagine, To doubt, To complain, To download, To upload}
\label{ch:ta-to-imagine-to-doubt-to-complain-to-download-to-upload}
\hlchaptermodality{230}

\begin{chapteropening}
\textbf{By the end of this chapter:} \emph{I can say to imagine, to doubt, to complain, to download and to upload.}

Complete the last lesson of chapter 230: I can say to imagine, to doubt, to complain, to download and to upload.
\end{chapteropening}

\section[\texorpdfstring{\ta{கற்பனை} \ta{செய்} (kaṟpaṉai cey)}{கற்பனை செய் (kaṟpaṉai cey)}]{\ta{கற்பனை} \ta{செய்} (kaṟpaṉai cey) — to imagine}

\label{lesson:TA-C230-karpanai-cey}



\begin{quote}
Before the new one: say the Tamil for to separate, then the Tamil for to throw.
\end{quote}

\subsection*{You'll want to know: \ta{கற்பனை} \ta{செய்}}
\textbf{\ta{கற்பனை} \ta{செய்}} — \emph{kaṟpaṉai cey} — \textquotedblleft{}to imagine\textquotedblright{}.

\begin{grammarlens}[title={What you've built}]
One more word for this chapter.
\end{grammarlens}

\subsection*{Your turn}
\begin{itemize}
\raggedright
  \item \emph{Say it:} \emph{kaṟpaṉai cey}
  \item \emph{Say it:} \emph{kaṟpaṉai cey}, once more
  \item \emph{Say it:} say \emph{kaṟpaṉai cey}
  \item \emph{Recall:} say the Tamil for to separate, then the Tamil for to throw, then say \emph{kaṟpaṉai cey} again
\end{itemize}

\begin{tcolorbox}[breakable,colback=teal!4,colframe=teal!35!black,title={Before you move on}]
What does \ta{கற்பனை} \ta{செய்} mean? (\textquotedblleft{}To imagine\textquotedblright{}.) Say it once more.
\end{tcolorbox}
\section[\texorpdfstring{\ta{சந்தேகப்படு} (cantēkappaṭu)}{சந்தேகப்படு (cantēkappaṭu)}]{\ta{சந்தேகப்படு} (cantēkappaṭu) — to doubt, to suspect}

\label{lesson:TA-C230-cantekappatu}



\begin{quote}
Before the new one: say the Tamil for to throw, then the Tamil for to imagine.
\end{quote}

\subsection*{You'll want to know: \ta{சந்தேகப்படு}}
\textbf{\ta{சந்தேகப்படு}} — \emph{cantēkappaṭu} — \textquotedblleft{}to doubt, to suspect\textquotedblright{}.

\begin{grammarlens}[title={What you've built}]
One more word for this chapter.
\end{grammarlens}

\subsection*{Your turn}
\begin{itemize}
\raggedright
  \item \emph{Say it:} \emph{cantēkappaṭu}
  \item \emph{Say it:} \emph{cantēkappaṭu}, once more
  \item \emph{Say it:} say \emph{cantēkappaṭu}
  \item \emph{Recall:} say the Tamil for to throw, then the Tamil for to imagine, then say \emph{cantēkappaṭu} again
\end{itemize}

\begin{tcolorbox}[breakable,colback=teal!4,colframe=teal!35!black,title={Before you move on}]
What does \ta{சந்தேகப்படு} mean? (\textquotedblleft{}To doubt, to suspect\textquotedblright{}.) Say it once more.
\end{tcolorbox}
\section[\texorpdfstring{\ta{குறை} \ta{சொல்} (kuṟai col)}{குறை சொல் (kuṟai col)}]{\ta{குறை} \ta{சொல்} (kuṟai col) — to complain, to find fault}

\label{lesson:TA-C230-kurai-col}



\begin{quote}
Before the new one: say the Tamil for to imagine, then the Tamil for to doubt.
\end{quote}

\subsection*{You'll want to know: \ta{குறை} \ta{சொல்}}
\textbf{\ta{குறை} \ta{சொல்}} — \emph{kuṟai col} — \textquotedblleft{}to complain, to find fault\textquotedblright{}.

\begin{grammarlens}[title={What you've built}]
One more word for this chapter.
\end{grammarlens}

\subsection*{Your turn}
\begin{itemize}
\raggedright
  \item \emph{Say it:} \emph{kuṟai col}
  \item \emph{Say it:} \emph{kuṟai col}, once more
  \item \emph{Say it:} say \emph{kuṟai col}
  \item \emph{Recall:} say the Tamil for to imagine, then the Tamil for to doubt, then say \emph{kuṟai col} again
\end{itemize}

\begin{tcolorbox}[breakable,colback=teal!4,colframe=teal!35!black,title={Before you move on}]
What does \ta{குறை} \ta{சொல்} mean? (\textquotedblleft{}To complain, to find fault\textquotedblright{}.) Say it once more.
\end{tcolorbox}
\section[\texorpdfstring{\ta{பதிவிறக்கு} (pativiṟakku)}{பதிவிறக்கு (pativiṟakku)}]{\ta{பதிவிறக்கு} (pativiṟakku) — to download}

\label{lesson:TA-C230-pativirakku}



\begin{quote}
Before the new one: say the Tamil for to doubt, then the Tamil for to complain.
\end{quote}

\subsection*{You'll want to know: \ta{பதிவிறக்கு}}
\textbf{\ta{பதிவிறக்கு}} — \emph{pativiṟakku} — \textquotedblleft{}to download\textquotedblright{}.

\begin{grammarlens}[title={What you've built}]
One more word for this chapter.
\end{grammarlens}

\subsection*{Your turn}
\begin{itemize}
\raggedright
  \item \emph{Say it:} \emph{pativiṟakku}
  \item \emph{Say it:} \emph{pativiṟakku}, once more
  \item \emph{Say it:} say \emph{pativiṟakku}
  \item \emph{Recall:} say the Tamil for to doubt, then the Tamil for to complain, then say \emph{pativiṟakku} again
\end{itemize}

\begin{tcolorbox}[breakable,colback=teal!4,colframe=teal!35!black,title={Before you move on}]
What does \ta{பதிவிறக்கு} mean? (\textquotedblleft{}To download\textquotedblright{}.) Say it once more.
\end{tcolorbox}
\section[\texorpdfstring{\ta{பதிவேற்று} (pativēṟṟu)}{பதிவேற்று (pativēṟṟu)}]{\ta{பதிவேற்று} (pativēṟṟu) — to upload}

\label{lesson:TA-C230-pativerru}



\begin{quote}
Before the new one: say the Tamil for to complain, then the Tamil for to download.
\end{quote}

\subsection*{You'll want to know: \ta{பதிவேற்று}}
\textbf{\ta{பதிவேற்று}} — \emph{pativēṟṟu} — \textquotedblleft{}to upload\textquotedblright{}.

\begin{grammarlens}[title={What you've built}]
One more word for this chapter.
\end{grammarlens}

\subsection*{Your turn}
\begin{itemize}
\raggedright
  \item \emph{Say it:} \emph{pativēṟṟu}
  \item \emph{Say it:} \emph{pativēṟṟu}, once more
  \item \emph{Say it:} say \emph{pativēṟṟu}
  \item \emph{Recall:} say the Tamil for to complain, then the Tamil for to download, then say \emph{pativēṟṟu} again
\end{itemize}

\begin{tcolorbox}[breakable,colback=teal!4,colframe=teal!35!black,title={Before you move on}]
What does \ta{பதிவேற்று} mean? (\textquotedblleft{}To upload\textquotedblright{}.) Say it once more.
\end{tcolorbox}
